\documentclass[11pt]{article}

\usepackage[margin=1in]{geometry}
\usepackage{amsmath,amssymb}
\usepackage{booktabs}
\usepackage{hyperref}
\usepackage{microtype}

\newcommand{\R}{\mathbb{R}}
\newcommand{\Xs}{\mathcal{X}}
\newcommand{\GP}{\mathcal{GP}}
\newcommand{\Norm}{\mathcal{N}}
\newcommand{\Data}{\mathcal{D}}
\newcommand{\Ulost}{U_{\text{lost}}}

\title{Separating Two Things Composite Modeling Provides:\\
Knowing the Outer Function from Observing the Intermediates}
\author{}
\date{}

\begin{document}
\maketitle

\begin{abstract}
\noindent
Composite Bayesian optimization gives a surrogate two advantages at once: it is
told the algebraic form of the outer function, and it gets to see the
intermediate quantities that function consumes. Every comparison in our
benchmark study changes both together, so none of them can say which advantage
produced the gain. This note designs, implements and validates a third
surrogate that receives exactly one of the two, which splits the composite gain
into its two parts. Part~\ref{part:overview} is a short summary of the whole
idea. The remaining parts give the full detail: the reasoning behind the
decomposition, the synthetic problems, the inference procedure for the new
surrogate and the three defects found in it, what is measured, and what has and
has not been verified, and the measured results.
\end{abstract}

\tableofcontents

% ======================================================================
\part{Overview}
\label{part:overview}
% ======================================================================

\section{The short version}

\paragraph{The setting.} Many expensive objectives factor as $f(x)=g(h(x))$,
where $h:\Xs\to\R^k$ is an expensive black box (a simulator, an experiment) and
$g:\R^k\to\R$ is a known, cheap function written down by whoever posed the
problem. A \emph{direct} surrogate fits a Gaussian process to $f$ and ignores
the factorization. A \emph{composite} surrogate fits Gaussian processes to $h$
and pushes their posterior samples through $g$. Our benchmark study found the
composite surrogate better on most problems, by margins running from $121\%$
down to nothing.

\paragraph{The problem with that comparison.} Going from direct to composite
changes two things simultaneously:
\begin{enumerate}
\item \textbf{knowledge of $g$} --- the surrogate is handed the outer function
and never has to learn it from data;
\item \textbf{observation of $h$} --- each expensive evaluation now returns $k$
numbers instead of one.
\end{enumerate}
Both are bundled into every result we have. So when the composite surrogate
wins by $121\%$ on DTLZ2 we cannot say which advantage did the work, and when it
wins by nothing on RGB we cannot say which one failed.

\paragraph{The fix.} Add a third surrogate that gets advantage (1) and not (2):
it knows $g$ exactly, holds the same prior belief about $h$ that the composite
surrogate holds, but only ever observes the scalar $y$. Call it the
\emph{latent} arm, because $h$ is never seen and must be inferred. Then
\begin{align}
\text{latent} - \text{direct} &= \text{the value of knowing } g,\\
\text{composite} - \text{latent} &= \text{the value of observing } h,
\end{align}
and the two sum exactly to the composite gain we already report.

\paragraph{What governs the split.} Observing $h$ can only help to the extent
that $y=g(h)$ fails to pin $h$ down. That makes invertibility of $g$ the
controlling property, and it yields two extreme cases we can check against
theory rather than intuition:
\begin{itemize}
\item \textbf{$g$ injective} $\Rightarrow$ $h=g^{-1}(y)$, so the latent and
composite arms see identical information and are the same model. Observing $h$
is worth \emph{exactly zero}.
\item \textbf{$g$ linear} $\Rightarrow$ a linear combination of independent GPs
is itself a GP, so the latent arm collapses into a direct GP on $y$. Knowing $g$
is worth \emph{approximately zero}. (Truss and RCM40, our two weakest results,
both have linear outer maps.)
\end{itemize}
Between those, the information $g$ destroys is the \emph{fiber} $g^{-1}(y)$: a
pair of points for $z^2$, several for $\cos$, a half-line for a clipped map, a
circle for a norm, a hyperbola for a product.

\paragraph{The experiment.} Nineteen outer maps applied to intermediates drawn
from the same GP family the surrogates assume, so the inner function is never a
confound. Single objective, $d=6$, $5$ initial points plus $40$ adaptive
evaluations, all three arms sharing the initial design and the acquisition
function so the surrogate is the only difference. Four independent draws of the
inner function, five initial designs each.

\paragraph{The hard part.} Arms 1 and 2 have closed-form posteriors --- ordinary
GP regression. The latent arm does not: it must infer quantities it never
observed through a nonlinear $g$, which is Markov chain Monte Carlo rather than
a formula. Every difficulty in this project has been there, and three genuine
defects were found and fixed (Section~\ref{sec:defects}).

\paragraph{Why we can believe the result.} Three of the nineteen conditions have
answers known in advance --- the identity map, an injective map, and a linear
map --- and are included solely so the machinery can be graded where grading is
possible. That is what caught the defects: on the identity map ``observing $h$''
must be exactly zero, and the code reported $+0.956$. After the fixes it
reports $-0.096\pm0.175$, and the two other controls pass as well.

\section{Reading guide}

Part~\ref{part:idea} gives the argument in full, including the worked cases
behind each prediction. Part~\ref{part:functions} specifies the synthetic
problems and explains the two that had to be redesigned. Part~\ref{part:latent}
is the inference procedure for the latent arm, the three defects and their
fixes. Part~\ref{part:measure} defines every reported quantity, including the
diagnostics. Part~\ref{part:protocol} is the protocol and the validation
evidence. Part~\ref{part:limits} states the limitations. The appendix maps each
piece of the design onto the code that implements it.

% ======================================================================
\part{The idea in detail}
\label{part:idea}
% ======================================================================

\section{The composite setting}

Fix a single objective $f(x)=g(h(x))$ on $\Xs=[0,1]^d$, with
$h:\Xs\to\R^k$ unknown and $g:\R^k\to\R$ known in closed form. After $n$
evaluations write $X=(x_1,\dots,x_n)$, $\mathbf{y}=(y_1,\dots,y_n)$ with
$y_i=f(x_i)+\varepsilon_i$, and
$\mathbf{h}(X)=(h(x_1),\dots,h(x_n))\in\R^{n\times k}$ for the matrix of
intermediates.

A surrogate's job at each iteration is to produce a predictive distribution over
$f(x')$ at candidate points $x'$, which an acquisition function turns into the
next evaluation. The three arms differ only in how that predictive is built.

\section{The confound, stated precisely}

Write $I_{\text{direct}}$, $I_{\text{comp}}$ for the information each surrogate
conditions on:
\begin{align}
I_{\text{direct}} &= \bigl(\mathbf{y}\bigr), \quad g \text{ unknown},\\
I_{\text{comp}}   &= \bigl(\mathbf{h}(X)\bigr), \quad g \text{ known}.
\end{align}
These differ in two coordinates at once. Any performance gap between them is a
sum of two effects with no way to separate them from the two numbers alone. The
latent arm supplies the missing corner of the $2\times2$ table:
\begin{center}
\begin{tabular}{lll}
\toprule
 & $g$ unknown & $g$ known \\
\midrule
observes $y$ only      & arm 1 (direct)  & arm 3 (latent) \\
observes $h$           & --- (incoherent) & arm 2 (composite) \\
\bottomrule
\end{tabular}
\end{center}
The empty cell is incoherent: a surrogate that observes $h$ but does not know
$g$ cannot predict $f$ at all. So three arms are the most that can be
instantiated, and three is enough --- the two available differences are exactly
the two effects, and they telescope.

\section{The three arms}

\begin{center}
\begin{tabular}{llll}
\toprule
Arm & Observes & Knows $g$ & Surrogate \\
\midrule
1. Direct    & $y_i=f(x_i)$ & no  & $f\sim\GP$ fit to $y$, hyperparameters by MLE \\
2. Composite & $h(x_i)$     & yes & $h_j\sim\GP$ conditioned on observed $h_j$, then apply $g$ \\
3. Latent    & $y_i=f(x_i)$ & yes & $h_j\sim\GP$ \emph{a priori}, $h$ inferred from $y$ \\
\bottomrule
\end{tabular}
\end{center}

\subsection{What the latent arm computes}

Its posterior over the unobserved intermediates is Bayes' rule and nothing more:
\begin{equation}
p\bigl(\mathbf{h}(X)\mid \mathbf{y}\bigr)\;\propto\;
\underbrace{\prod_{i=1}^{n}\Norm\!\bigl(y_i;\,g(h(x_i)),\,\sigma^2\bigr)}_{\text{likelihood: does this }h\text{ explain }y?}
\;\cdot\;
\underbrace{\prod_{j=1}^{k}\Norm\!\bigl(h_j(X);\,m_j,\,v_jK\bigr)}_{\text{GP prior: is this }h\text{ a plausible function?}} ,
\label{eq:posterior}
\end{equation}
where $K_{ii'}=k_{\mathrm{RBF}}(x_i,x_{i'})$ and $(m_j,v_j)$ are the prior mean
and variance of the $j$th intermediate.

Prediction is then exact \emph{given} the latents, because
$h(x')\mid\mathbf{h}(X)$ is an ordinary GP conditional. So the arm-3 predictive
is a \textbf{mixture of Gaussian processes}, one component per posterior sample
of $\mathbf{h}(X)$:
\begin{equation}
p\bigl(f(x')\mid\mathbf{y}\bigr)=\int p\bigl(g(h(x'))\mid \mathbf{h}(X)\bigr)\,
p\bigl(\mathbf{h}(X)\mid\mathbf{y}\bigr)\,d\mathbf{h}(X)
\;\approx\;\frac{1}{S}\sum_{s=1}^{S}p\bigl(g(h(x'))\mid\mathbf{h}^{(s)}\bigr).
\end{equation}
That mixture is the whole point. A direct GP on $y$ is a single Gaussian and
cannot represent multimodality; if two very different intermediate functions
both explain the data, the latent arm carries both hypotheses forward and its
predictive is genuinely bimodal.

\subsection{The prior, and what ``sharing'' it means}

A GP prior here is three numbers per intermediate: a mean $m_j$, a variance
$v_j$, and a lengthscale $\ell$ shared across coordinates. Everything else the
GP does follows from those plus the data.

Arms 2 and 3 are constructed from the \textbf{same} prior object --- same
$(m_j,v_j,\ell)$, same conditioning code --- and it is not an estimate: it is
the recipe the inner functions were actually generated from
(Section~\ref{sec:draws}). The only difference between the two arms is therefore
what gets conditioned on: arm 2 conditions on $\mathbf{h}(X)$ directly, arm 3
must first infer it from $\mathbf{y}$. That is deliberate. Had arm 3 been given
a different lengthscale or variance, the gap between arms 2 and 3 would blend
``better prior'' with ``better observations'' and measure neither.

Arm 1 cannot share it and does not: it models $y$, and for a general $g$ the
function $g(h(\cdot))$ is not a GP, so there is no true prior to hand it. It
estimates its own by maximum likelihood. Section~\ref{sec:hyperconfound} deals
with the asymmetry this creates.

\subsection{The accounting identity}

With $R_a$ the performance of arm $a$ on a common metric,
\begin{equation}
\underbrace{(R_1-R_3)}_{\text{knowing }g}\;+\;\underbrace{(R_3-R_2)}_{\text{observing }h}
\;=\;\underbrace{(R_1-R_2)}_{\text{composite gain we already report}} .
\end{equation}
The identity is trivially true, which is a feature: the right-hand side is a
number we have already published for each benchmark, so the decomposition is
guaranteed to add up and any failure to do so is a bug, not a finding.

\section{Why the decomposition should be informative}

\subsection{Invertible $g$: observing $h$ is worth nothing}

If $g$ is injective and observations are noise-free, $h(x_i)=g^{-1}(y_i)$
exactly. Arm 3 recovers the intermediates from the objective alone, so arms 2
and 3 condition on identical information and are literally the same model. Any
composite advantage on such a problem is attributable entirely to knowing $g$.

The exponential warping family $g(z)=\exp(\beta z)$ is of this type. Knowing $g$
is still worth something --- it tells the surrogate that $y>0$, and how the scale
is warped --- but the extra numbers are redundant.

Two of these maps are in the run specifically as a sanity check: if the code
reports anything other than exactly zero for ``observing $h$'' on $\exp$ or
$\sinh$, it is broken.

\subsection{Linear $g$: knowing $g$ is worth nothing}

If $g(\mathbf{y})=\sum_j c_jy_j$ and the $h_j$ are independent GPs, then
$g(h(\cdot))$ is \emph{itself} a GP with mean $\sum_jc_jm_j$ and kernel
$\bigl(\sum_jc_j^2v_j\bigr)k(\cdot,\cdot)$. Arm 3 therefore collapses to an
ordinary GP fitted to $y$ --- which is arm 1 with a particular kernel. Up to
hyperparameter estimation, knowing $g$ buys nothing.

This is the sharpest prediction in the design, and it bears directly on the
paper: Truss and RCM40 both have linear outer maps and both produced our weakest
results. If the prediction holds, the explanation is not the vague ``linear maps
are easy'' but something specific --- on a linear map there is no knowledge to
exploit, so whatever gain the composite surrogate shows must come from the extra
observations alone.

\subsection{Non-invertible $g$: both channels are live}

Take the smallest interesting case, $g(z)=z^2$, noise-free. Then $h(x)$ is known
only up to sign:
\begin{equation}
p\bigl(h(x')\mid g(h(x))=y\bigr)
=p\bigl(h(x')\;\bigm|\;h(x)\in\{+\sqrt{y},\,-\sqrt{y}\}\bigr),
\label{eq:square}
\end{equation}
so the posterior is a two-component mixture after one observation, and a mixture
over $2^n$ sign patterns after $n$, each pattern weighted by its plausibility
under the GP prior. Crucially the weights are \emph{not} uniform: a pattern that
requires the function to jump between $+\sqrt{y}$ and $-\sqrt{y}$ at two nearby
inputs is heavily penalized by the prior, so the data do constrain the signs
jointly even though each observation alone does not. Observing $h$ removes the
ambiguity outright, which is what arm 2 measures.

But knowing $g$ is worth something even with no observation of $h$ at all. The
range of $g$ is $[0,\infty)$, so arm 3 assigns
\begin{equation}
P\bigl(f(x')=-1\mid\Data\bigr)=P\bigl(h(x')^2=-1\mid\Data\bigr)=0,
\end{equation}
whereas a direct GP on $y$ --- a Gaussian --- puts positive probability on values
that can never occur. More generally, knowing $g$ imposes the support, the
curvature, and any kinks or saturation of the outer map, none of which a
stationary GP on $y$ can represent.

In general the information destroyed by $g$ is the \emph{fiber} $g^{-1}(y)$: a
finite set for $z^2$ or $\cos$, a half-line for a clipped or rectified map, a
circle for a Euclidean norm, a hyperbola for a product. This is the structure
the synthetic maps in Part~\ref{part:functions} are chosen to vary.

\subsection{Predictions, stated before the run}

\begin{center}
\begin{tabular}{llcc}
\toprule
$g$ & Type & knowing $g$ & observing $h$ \\
& & (arm 1 $-$ arm 3) & (arm 3 $-$ arm 2) \\
\midrule
$z$                               & identity       & $0$ & $0$ \\
$\exp(\beta z)$, $\sinh(\beta z)$ & invertible     & $+$ & $0$ by construction \\
$(z-c)^2$                         & lossy, tunable & $+$ & $+$ at $c=0$, $\to0$ as $c$ grows \\
$|z|$                             & lossy          & $+$ & $+$ \\
$y_1+y_2$                         & linear         & $\approx0$ & $+$ \\
norm, max, product, ratio         & lossy          & $+$ & $+$ \\
$-\max(z,0)$                      & censors away from optimum & $+$ & $\approx 0$ \\
$\mathrm{clip}(2y_1,-1,1)+y_2$    & censors at the optimum & $+$ & $+$ \\
\bottomrule
\end{tabular}
\end{center}

The last two rows deserve emphasis because they are the design's one genuinely
non-obvious prediction. Both maps destroy information by censoring, but
$-\max(z,0)$ is optimized at the largest $h$, where $g$ is injective, so the
destroyed information sits \emph{away from} the optimum; the clipped pair is
constructed so the censored coordinate matters at the optimum. If observing $h$
pays off on the second and not the first, then what matters is not how much
information $g$ destroys but \emph{where}.

% ======================================================================
\part{The test functions}
\label{part:functions}
% ======================================================================

\section{Design principle: the inner function must not be a confound}
\label{sec:draws}

If the inner functions were arbitrary analytic expressions, arms 2 and 3 might
win or lose because a GP happens to model those expressions well or badly, and
we would learn nothing about $g$. So every primitive $h_j$ is a fixed draw from
the \emph{same} RBF-GP family the surrogates assume, generated with random
Fourier features:
\begin{equation}
h(x)=\sqrt{\tfrac{2}{M}}\sum_{m=1}^{M}a_m\cos(\omega_m^\top x+b_m),
\qquad a_m\sim\Norm(0,1),\;b_m\sim U(0,2\pi),\;\omega_m\sim\Norm(0,\ell^{-2}I),
\end{equation}
with $M=1024$ and $\ell=0.5$. As $M\to\infty$ this converges to a draw from
$\GP(0,k_{\mathrm{RBF}})$; at $M=1024$ it is close enough that the prior handed
to arms 2 and 3 is correct rather than approximately correct.

The draw is standardized on a fixed Sobol reference set, so the prior on the
standardized function is $\GP(-\mu/s,\,k_{\mathrm{RBF}}/s^2)$ with $\mu,s$ the
measured mean and standard deviation. Those two constants are exactly what arms
2 and 3 receive as $(m_j,v_j)$. The same draw is reused across every outer map,
so switching $g$ never changes the underlying function.

\paragraph{Why standardizing matters beyond scale.} It puts the draw's mean at
zero, so $h$ takes \emph{both signs} over the domain. Without that, $z^2$ would
be injective on the region actually visited and would destroy nothing --- the
sweep's whole knob would be dead.

\paragraph{Objective normalization.} Each final objective is mapped to roughly
$[0,1]$ using the $1$st and $99$th percentiles of $g(h(\cdot))$ on a fixed
reference set, so outer maps do not differ merely in scale. The transform is
affine and known, so it is folded into $g$ for all three arms and changes
nothing about who knows what.

\section{Function draws versus seeds}
\label{sec:seeds}

Two independent sources of randomness must be varied, and they are not
interchangeable.

\begin{description}
\item[The initial design.] Which $5$ Sobol points start the run. All three arms
share it within a seed, which is what makes the comparison paired.
\item[The function draw.] Which random $h$ was drawn in the first place.
\end{description}

Only the first was varied in the initial version of this study: the
function-draw seed was fixed, so twenty ``seeds'' meant twenty initial designs
on \emph{one} sampled function, and the reported standard errors described
variation in the starting points only. A conclusion could have been a property
of that one draw with no way to tell from the numbers. The run now sweeps four
draws with five designs each --- identical total compute --- and the collection
step pools them, so the error bars cover both sources.

\section{The maps}

Nineteen conditions in four groups.

\paragraph{Controls (3).} $g(z)=z$; $g(\mathbf{y})=y_1+y_2$; $g(z)=\exp(z)$.
Identity checks the machinery end to end (both effects are zero). The linear map
checks the sampler against an exactly computable posterior. The injective map
checks that ``observing $h$'' comes out at zero when theory says it must.

\paragraph{The sweep (5).} $g_c(z)=(z-c)^2$ for $c\in\{0,0.5,1,2,3\}$ holds the
functional form fixed and varies only how much ambiguity there is. At $c=0$ the
sign of $h$ is lost wherever the draw crosses zero; at $c=3$ the shifted argument
almost never changes sign over the domain, so $g$ is effectively injective. This
is the cleanest available knob on ``how much $g$ discards,'' and the central plot
of the study is the ``observing $h$'' gain against the measured $\Ulost$ along
it.

\paragraph{Matched pairs (9).} Each pair changes one property and holds the rest:
\begin{center}
\begin{tabular}{lll}
\toprule
Pair & Shared & Isolates \\
\midrule
$z^2$ vs.\ $|z|$ & both two-to-one & smoothness at the optimum \\
$\exp(\beta z)$ vs.\ $\sinh(\beta z)$ & both injective, both smooth & whether a range constraint helps \\
$\cos(\omega z)$, $\omega\in\{1,2,4\}$ & smooth, periodic & \emph{how many} preimages \\
$-\max(z,0)$ vs.\ clipped pair & both censor & \emph{where} the censoring bites \\
\bottomrule
\end{tabular}
\end{center}

\paragraph{Benchmark analogs (5).} Two intermediates into one objective, each
mirroring a real benchmark: $y_1+y_2$ (Truss, RCM40), $y_1y_2$ (DTLZ2),
$\sqrt{y_1^2+y_2^2}$ (Color), $\max(y_1,y_2)$ (Lake), and
$u_1/(0.05+u_1+\lambda u_2)$ with $u=\mathrm{sigmoid}(y)$ (RGB). These let the
synthetic conclusions be carried back to the benchmark results.

\section{Two maps that had to be redesigned}
\label{sec:degenerate}

The first version of the censoring pair was degenerate, and the failure is
instructive.

\paragraph{The symptom.} On $g(z)=\max(z,0)$ and
$g(z)=\mathrm{clip}(cz,-1,1)$, all three arms reported a final $\log_{10}$
regret of exactly $-8.000$, which is the clamp floor in the metric --- i.e.\ zero
regret for every arm on every seed.

\paragraph{The cause.} Both maps are minimized on a flat region of positive
measure. A bare $\max(z,0)$ takes its minimum value $0$ on the entire censored
half $\{h\le0\}$, and a bare clip takes its minimum on the whole lower
saturation region. The five-point Sobol initial design already contains such a
point with high probability, so every arm starts at the optimum and the
condition carries no signal about anything.

\paragraph{The fix.} Both are now minimized at a unique point while censoring the
same set:
\begin{align}
g(z)&=-\max(z,0) && \text{optimum at }\max h\text{, unique; censors }\{h\le0\}\\
g(\mathbf{y})&=\mathrm{clip}(2y_1,-1,1)+y_2 && \text{optimum is a genuine tradeoff}
\end{align}
The negation puts the optimum where $g$ is injective, which is what turns this
pair into the ``where does the censoring bite'' contrast of the previous section.
The clipped version adds a second, uncensored intermediate so the objective is
not flat, while the first coordinate stays censored inside saturation.

\paragraph{Verification.} The median final regret of a five-point Sobol design
alone rose from $0$ to $1.02$ for the rectified map and $0.62$ for the clipped
pair, i.e.\ there is now something to optimize. A four-seed pilot separated the
arms as predicted: $\Ulost=0.088$ with ``observing $h$'' at $-0.045\pm0.194$ for
the rectified map, against $\Ulost=0.459$ and $+2.811\pm1.295$ for the clipped
pair.

% ======================================================================
\part{Inference for the latent arm}
\label{part:latent}
% ======================================================================

\section{Why this needs a sampler at all}

Arms 1 and 2 are ordinary GP regression: the posterior is Gaussian and available
in closed form. Arm 3 is not, and the reason is structural rather than technical.
Its likelihood in \eqref{eq:posterior} depends on the latent $\mathbf{h}(X)$
through $g$, which is nonlinear; a Gaussian prior times a non-Gaussian
likelihood is not Gaussian, so there is no formula. Asking ``what would knowing
$g$ alone be worth?'' necessarily means asking a surrogate to reason about
quantities it never observed, and that is an inference problem the other two arms
simply do not have.

The state being sampled is small --- $n\times k$ numbers, at most $45\times2=90$
at our budget --- but it is awkward: the posterior is multimodal by construction
(that is the phenomenon under study), the likelihood is sharply peaked when
$\sigma$ is small, and several of the maps have kinks or flat regions so
gradients are unavailable or uninformative.

\section{Closed forms where they exist}

\begin{description}
\item[Identity and injective maps.] $h=g^{-1}(y)$ exactly, so arm 3 is
implemented as arm 2. This is a shortcut, not a different model.
\item[Linear maps.] The posterior \eqref{eq:posterior} is Gaussian and available
analytically; arm 3 samples it directly.
\end{description}
Both shortcuts can be disabled with a flag so the general sampler runs on those
maps too, which is how they are validated (Section~\ref{sec:validation}).

\section{The general sampler}

Two move types alternate.

\subsection{Elliptical slice sampling}

Elliptical slice sampling \cite{murray2010} is the standard sampler for a GP
prior with an arbitrary likelihood. Given the current state $\mathbf{h}$ it draws
$\boldsymbol{\nu}$ from the prior, proposes
\begin{equation}
\mathbf{h}(\theta)=\mathbf{h}\cos\theta+\boldsymbol{\nu}\sin\theta,
\end{equation}
and shrinks the bracket on $\theta$ until the likelihood threshold is met. The
ellipse is constructed so that every point on it has the same prior density,
which is why no acceptance ratio is needed. It requires no gradients ---
essential, because $|z|$, $\max$ and the clipped map are non-differentiable ---
and no step size to tune.

\subsection{Level-set moves, and why they are necessary}

Elliptical slice sampling alone is not enough. When $\sigma$ is small the
likelihood is sharply peaked on the constraint set
$\{\mathbf{h}: g(h(x_i))=y_i\ \forall i\}$, so a proposal must shrink to
$\theta\sim\sigma/\mathrm{sd}(h)$ to remain inside it. Movement \emph{along} the
constraint set is then a very slow random walk --- and that set is precisely the
ambiguity the experiment is about. A sampler that cannot traverse it will
confidently report one branch and miss the multimodality entirely.

Those directions are known in closed form for each $g$. We therefore propose
$h(x_i)\mapsto h'(x_i)$ with
\begin{equation}
g\bigl(h'(x_i)\bigr)=g\bigl(h(x_i)\bigr)\quad\text{exactly},
\end{equation}
from a symmetric proposal density. The likelihood ratio is then identically one
and the Metropolis--Hastings acceptance probability reduces to the GP prior ratio
alone. With $Q=K^{-1}$ and $r_j=h_j(X)-m_j$, changing coordinate $i$ of latent
$j$ from $\alpha$ to $\beta$ gives
\begin{equation}
\Delta\log p=-\frac{1}{2v_j}\Bigl[Q_{ii}\bigl(\beta^2-\alpha^2\bigr)
+2(\beta-\alpha)\bigl((Qr_j)_i-Q_{ii}\alpha\bigr)\Bigr],
\label{eq:priorratio}
\end{equation}
which is $O(n)$ per single-site move given $Qr_j$. Sites are visited in random
scan order.

The moves are the symmetries of each map:
\begin{center}
\begin{tabular}{ll}
\toprule
Map & Level-set move \\
\midrule
$(z-c)^2$, $|z|$        & reflection $h\mapsto 2c-h$ \\
$\cos(\omega z)$        & reflection through the nearest multiple of $\pi/\omega$; period shift \\
$-\max(z,0)$            & random walk confined to the censored region \\
clipped pair            & random walk on the clipped coordinate within one saturation region \\
$\sqrt{y_1^2+y_2^2}$    & rotation of $(h_1,h_2)$ \\
$y_1y_2$                & negation of both; reciprocal rescaling $(ch_1,h_2/c)$ \\
$\max(y_1,y_2)$         & coordinate swap; move the smaller coordinate \\
$y_1+y_2$               & shift $(h_1+\delta,h_2-\delta)$ \\
\bottomrule
\end{tabular}
\end{center}
Only the ratio map has no convenient exact symmetry and relies on elliptical
slice sampling alone --- noted as a limitation in Part~\ref{part:limits}.

\section{Three defects, and what they teach}
\label{sec:defects}

The first full run produced a table that looked entirely reasonable: smooth
trends across the sweep, sensible signs everywhere, tight standard errors. It
was wrong. The identity control read ``observing $h$'' $=+0.956\pm0.274$ where
the true value is exactly $0$ by construction. Three separate defects
contributed.

\subsection{Defect 1: collapsed Monte Carlo pairing}

The predictive samples were formed by giving every latent sample a
\emph{shared} noise draw. When the latent posterior is tight, samples then
differ only through their nearly identical posterior means, so the effective
sample size collapses to one. Fixed by drawing $n$ independent noise vectors and
pairing each with a latent sample cycled from the chain, matching the composite
arm's sample count so the two are comparable. This was a real bug but not the
main one: the identity control barely moved.

\subsection{Defect 2: a false lead}

The next hypothesis was that $\sigma$ was acting as a resolution floor. It was
tested directly at $\sigma\in\{10^{-2},10^{-3},10^{-4}\}$, which gave
$+1.036$, $+0.883$, $+1.156$ --- no trend --- and $\Ulost$ moved only from
$0.0026$ to $0.0006$ when it should scale as $\sigma^2$. That non-scaling was
itself the clue: something other than the likelihood was setting the spread.

\subsection{Defect 3: latents never reaching the constraint surface}

Two compounding mistakes in the warm-start path.

\begin{enumerate}
\item Burn-in anneals $\sigma$ downward so a chain started from the prior can
find the data. That annealing was also being applied on warm starts, re-heating
latents that were already correct; the short warm burn-in could not recover them.
\item A newly appended BO point was initialized at its GP conditional mean,
which is \textbf{not} on the surface $\{g(h)=y\}$. As established above,
elliptical slice sampling cannot travel to that surface once the likelihood is
tight --- its steps are $O(\sigma/\mathrm{sd}(h))$.
\end{enumerate}

So each BO iteration appended a point that stayed off the surface, and the error
accumulated. Measured constraint residual across successive iterations of one
run:
\begin{center}
\begin{tabular}{lrrrrrr}
\toprule
BO iteration & 1 & 2 & 3 & 4 & 5 & 6 \\
\midrule
mean $\|g(h)-y\|$ & $0.00013$ & $0.00435$ & $0.02666$ & $0.01903$ & $0.04255$ & $0.12144$ \\
\bottomrule
\end{tabular}
\end{center}

The fix is in two parts: do not anneal on a warm start, and \emph{project} new
rows onto the constraint surface before sampling, by local gradient descent on
$\|g(h)-y\|^2$ started from the GP conditional mean. Descent lands on the
preimage nearest that mean, which keeps whichever branch the prior prefers; the
level-set moves then explore the others. Projection is used only to initialize,
so it biases nothing --- the stationary distribution is still
\eqref{eq:posterior}.

After the fix the identity control reads $-0.096\pm0.175$, $\Ulost=0$ exactly,
and the residual is stable at $\approx0.0015$ across iterations.

\subsection{Why this matters methodologically}

None of the three defects was visible in the results. Each produced plausible
numbers with plausible error bars, and the trends across the sweep were in the
predicted direction throughout. What exposed them was a condition whose answer
was known before the run. A study containing only the interesting maps would have
produced a confident, wrong table --- and the wrongness pointed in the direction
of the hypothesis, which is the worst case.

The practical rule the design now follows: every arm is graded on at least one
condition where the correct answer is available independently of the code being
graded.

% ======================================================================
\part{What is measured}
\label{part:measure}
% ======================================================================

\section{Information destroyed by $g$}

\begin{equation}
\Ulost=\frac{1}{nk}\sum_{i=1}^{n}\sum_{j=1}^{k}
\operatorname{Var}\bigl(h_j(x_i)\mid \mathbf{y}\bigr).
\end{equation}
This is the posterior variance of the intermediates \emph{at the points already
evaluated} --- not at new points, where uncertainty would reflect ordinary GP
extrapolation rather than anything about $g$. It is zero when $g$ is injective
and noise-free, and grows with the size of the fiber. It is the measured version
of the quantity the whole design is organized around, and it is what the
``observing $h$'' gain is plotted against.

\section{Optimization performance}

Each arm reports simple regret against a reference minimum $f^\star$; the
headline number is the final
\begin{equation}
\log_{10}\bigl(\max(f_{\text{best}}-f^\star,\,10^{-8})\bigr).
\end{equation}
The log is used because regret spans orders of magnitude across maps, and the
clamp keeps a run that finds the optimum exactly from producing $-\infty$. It
also means a reported $-8.000$ is a saturated value, which is how the degenerate
maps of Section~\ref{sec:degenerate} were spotted.

Comparisons are paired: all three arms share an initial design within a seed, so
the reported difference is the mean of the per-seed differences and its standard
error is
\begin{equation}
\mathrm{SE}=\sqrt{\frac{1}{S(S-1)}\sum_{s=1}^{S}\bigl(d_s-\bar d\bigr)^2},
\qquad d_s=R_a^{(s)}-R_b^{(s)} .
\end{equation}
Pairing matters a great deal here: the variation between initial designs is
larger than the effects being measured, and differencing within a seed removes
it.

\section{Sampler diagnostics, and what each one means}

Three numbers are reported per condition. They answer different questions and
the third is the one to trust.

\begin{description}
\item[$\hat R$ (Gelman--Rubin).] Several chains are started independently;
$\hat R$ compares between-chain spread to within-chain spread, and equals $1$
when they are indistinguishable. The conventional threshold is $1.1$. It is
\textbf{unreliable here} whenever the posterior is nearly a point, because it
divides by within-chain variance --- which is exactly the injective and
large-$c$ cases.

\item[Chain disagreement.] The largest gap between any two chains' predicted
objective value at held-out points, in normalized objective units where the
range is about $1$. This is reported because it is the quantity BO actually
cares about, and because it stays meaningful when $\hat R$ does not. Under
$\approx0.02$ is comfortable.

\item[Move acceptance.] The fraction of proposed level-set moves accepted. Near
zero means the sampler is confined to one branch --- which is a \emph{failure} if
other branches are plausible and \emph{correct} if they are not. It should
therefore track $\Ulost$, and on the sweep it does: both fall to zero together.
\end{description}

Critically, all diagnostics are computed on the predictive of $f$, not on
$\mathbf{h}$ itself. Under an even $g$ a chain that has globally flipped the sign
of $h$ represents the identical function; measuring on $h$ would register that as
catastrophic non-convergence when nothing is wrong.

\section{A confound that is not yet resolved}
\label{sec:hyperconfound}

Arms 2 and 3 are given the generator's true kernel. Arm 1 must estimate its
hyperparameters by maximum likelihood, because for a general $g$ there is no
true kernel for $f$. The ``knowing $g$'' comparison therefore mixes the value of
knowing $g$ with the value of having a correct prior instead of one fitted to
between $5$ and $45$ points.

The identity map measures the second effect \emph{alone}, since knowing $g$ is
worth exactly nothing there. In the pre-fix run it read $+0.628\pm0.260$ --- the
same order of magnitude as most of the real effects, so this is not a minor
correction. Two ways to handle it:
\begin{enumerate}
\item \textbf{Subtract the baseline}, assuming the hyperparameter advantage is
roughly constant across maps. Cheap, and available from the current run, but the
assumption is unverified.
\item \textbf{Re-run with maximum-likelihood hyperparameters in all three
arms}, which removes the asymmetry rather than correcting for it. This is the
honest version and is the main outstanding piece of work.
\end{enumerate}
Until (2) is done, ``knowing $g$'' should be read as an upper bound.

% ======================================================================
\part{Protocol and validation}
\label{part:protocol}
% ======================================================================

\section{Protocol}

Single objective on $[0,1]^6$ with $\ell=0.5$; $n_0=5$ Sobol points and $40$
adaptive evaluations, matching Experiment~1 of the benchmark study; four
function draws $\times$ five initial designs per condition, pooled. All three
arms share the initial design, the acquisition function, the candidate pool and
the budget, so the surrogate is the only thing that differs.

The acquisition is Monte Carlo LogEI in a numerically stable form,
\begin{equation}
\log\mathrm{EI}(x')\approx\log\!\Bigl[\tfrac{1}{N}\sum_{s}\tau\,
\mathrm{softplus}\bigl((f_{\text{best}}-f_s(x'))/\tau\bigr)\Bigr],
\end{equation}
evaluated by \texttt{logsumexp} over $N=256$ samples with $\tau=10^{-3}$, on a
$2048$-point Sobol pool with local refinement of the best four candidates. The
softplus smoothing matters for arm 3 specifically: with a mixture predictive the
hard hinge $\max(0,\cdot)$ makes the acquisition surface non-smooth in a way that
interacts badly with local refinement.

Runtime: about $2.5$ minutes per (map, draw) single-threaded, so the full
$19\times4$ grid is roughly $20$--$30$ minutes on $16$ workers. Single-threaded
is deliberate --- the matrices are at most $45\times45$ and BLAS threading
measurably does not help at that size, so many one-thread jobs beat few
many-thread ones.

\section{Validation}
\label{sec:validation}

Three checks, each on a condition whose answer is available independently. They
are the reason the defects of Section~\ref{sec:defects} were caught, and they are
run before every experiment.

\subsection{Check 1: against an exactly computable posterior}

On the linear map with $n=20$ and the shortcut disabled, the posterior is
Gaussian in closed form, so the sampler can be graded directly.

\begin{center}
\begin{tabular}{lll}
\toprule
Quantity & Measured & Reading \\
\midrule
$|$mean $-$ exact mean$|$ / exact sd & median $0.069$, max $0.364$ & off by $7\%$ of a posterior sd \\
sampled sd / exact sd & median $0.968$, range $[0.84,1.09]$ & $3\%$ too narrow on average \\
move acceptance & $0.79$ & branches are being explored \\
$\hat R$ / chain disagreement & $1.150$ / $0.0085$ & chains agree to $0.85\%$ of the range \\
\bottomrule
\end{tabular}
\end{center}

The second row is the one that matters most: too-narrow posteriors make the
surrogate overconfident, which stops it exploring. $\hat R$ exceeds $1.1$ here,
which is the case discussed above where it is not informative --- the
disagreement figure is the reliable one.

\subsection{Check 2: calibration against known truth}

On $\exp(z)$ with the shortcut disabled, the true $h$ is recoverable, so we can
ask whether the sampler's error bars contain it. If calibrated,
$z=|$posterior mean $-$ true $h|/$posterior sd behaves like $|\Norm(0,1)|$:
median $\approx0.67$, about $95\%$ below $2$.

Measured: median $0.19$, max $1.10$, $\mathbf{100\%}$ within two standard
deviations. The posteriors are therefore somewhat \emph{wider} than necessary,
which is the safe direction. Before the fixes this read $95\%$ with a largest raw
error of $0.620$; it is now $100\%$ with $0.386$.

The largest remaining error sits at $h=-1.96$, where $\exp(h)\approx0.14$ and the
slope of $g$ falls below the noise level $\sigma=10^{-2}$. This is correct
behavior and a substantive caveat: \textbf{with observation noise, an injective
$g$ still destroys information wherever it is flat}, so ``arm 2 equals arm 3 for
injective $g$'' holds exactly only in the noise-free limit. The same map reports
$\Ulost=2.4\times10^{-2}$ for this reason, rather than $0$.

\subsection{Check 3: does the designed knob actually turn?}

\begin{center}
\begin{tabular}{lrrrr}
\toprule
Map & $\Ulost$ & $\hat R$ & Chain disagreement & Move acceptance \\
\midrule
$(z-0)^2$   & $0.8703$ & $1.013$ & $0.0272$ & $0.68$ \\
$(z-0.5)^2$ & $1.0789$ & $1.014$ & $0.0331$ & $0.54$ \\
$(z-1)^2$   & $0.5154$ & $1.129$ & $0.0356$ & $0.32$ \\
$(z-2)^2$   & $0.1075$ & $1.090$ & $0.0261$ & $0.07$ \\
$(z-3)^2$   & $0.0038$ & $1.226$ & $0.0135$ & $0.00$ \\
$|z|$       & $0.8812$ & $1.009$ & $0.0356$ & $0.66$ \\
\bottomrule
\end{tabular}
\end{center}

Information destroyed falls by more than two orders of magnitude across the
sweep, and move acceptance falls with it: by $c=3$ the reflected proposal is so
implausible under the prior that it is never accepted, which is the correct
behavior when no ambiguity remains. Two independent measurements of the same
collapse.

\paragraph{One blemish, recorded honestly.} $c=0.5$ reads \emph{higher} than
$c=0$ ($1.079$ vs.\ $0.870$), so $\Ulost$ is not monotone along the sweep. The
automated check only tests the endpoints, so it passes. The likely explanation is
that at $c=0.5$ both preimages are still live but now asymmetric, giving a wide
lopsided two-bump posterior whose variance genuinely exceeds that of the
symmetric case; variance is not a monotone function of ``how much ambiguity''
when the modes are unequally weighted. Whether it survives pooling over function
draws is an open question, and if it does, $\Ulost$ should be supplemented by a
mode-count or entropy measure.

\paragraph{Chain disagreement on the sweep.} It reads $0.026$--$0.036$, above the
$0.02$ guideline, on exactly the maps whose posteriors are genuinely multimodal.
This is not a defect --- it is real posterior mass being reported --- but it means
the sweep's numbers carry extra Monte Carlo noise on top of the seed variation.
Pooling over function draws mitigates it, since draws average over which modes
happen to matter.

\section{Results}
\label{sec:results}

Nineteen conditions, four function draws, five initial designs each: $380$ BO
runs per arm. Everything below is pooled over the twenty (draw, seed) pairs with
the paired standard error of Section~\ref{part:measure}.

\subsection{How to read the numbers}

The reported quantity is a difference of final $\log_{10}$ regrets, so it
converts to a multiplicative factor: an ``observing $h$'' value of $d$ means the
composite arm ends with $10^{d}$ times lower regret than the latent arm. A value
of $1$ is a tenfold improvement; a value of $0$ means the two arms are
indistinguishable.

\paragraph{The bias correction.} On \texttt{identity} with the shortcut
disabled --- so the general sampler runs on a problem where $h$ is exactly
recoverable and $\Ulost=0.0016$ --- ``observing $h$'' measures
$+0.409\pm0.133$ where the true value is $0$. The latent arm therefore carries a
handicap unrelated to information, discussed in
Section~\ref{sec:residual}. Every map whose latent arm invokes the sampler is
reported both raw and with $0.409$ subtracted. The three maps that use the
closed-form shortcut need no correction: there the latent arm \emph{is} the
composite arm and the measurement is exact.

\subsection{What each condition tests}

\begin{center}
\small
\begin{tabular}{llp{3.6cm}p{3.3cm}}
\toprule
Map & $g$ & What $g$ destroys & Hypothesis \\
\midrule
\multicolumn{4}{l}{\emph{Controls --- answers known before the run}}\\
\texttt{identity}   & $z$ & nothing & both effects exactly $0$ \\
\texttt{exp\_b1}    & $e^{z}$ & nothing (monotone, positive) & observing $h$ exactly $0$ \\
\texttt{sinh\_b1}   & $\sinh z$ & nothing (monotone, unbounded) & observing $h$ exactly $0$ \\
\texttt{sum}        & $y_1+y_2$ & how the sum splits & knowing $g$ $\approx0$ \\
\midrule
\multicolumn{4}{l}{\emph{The sweep --- one form, information loss tuned by $c$}}\\
\texttt{shiftsq\_c0}   & $(z-0)^2$   & sign of $z$   & largest gain \\
\texttt{shiftsq\_c0.5} & $(z-0.5)^2$ & sign of $z-0.5$ & \\
\texttt{shiftsq\_c1}   & $(z-1)^2$   & sign of $z-1$ & gain falls \\
\texttt{shiftsq\_c2}   & $(z-2)^2$   & sign of $z-2$ (rarely ambiguous) & with $c$ \\
\texttt{shiftsq\_c3}   & $(z-3)^2$   & almost nothing & gain $\to0$ \\
\midrule
\multicolumn{4}{l}{\emph{Matched pairs --- one property changed at a time}}\\
\texttt{abs}     & $|z|$ & sign of $z$, with a kink & as \texttt{shiftsq\_c0}; kink is the only difference \\
\texttt{cos\_w1} & $\cos z$    & which of several preimages & gain grows with the \\
\texttt{cos\_w2} & $\cos 2z$   & which of more preimages & number of preimages \\
\texttt{cos\_w4} & $\cos 4z$   & which of many preimages & \\
\texttt{relu}    & $-\max(z,0)$ & everything below $0$, \emph{away from} the optimum & gain $\approx0$ \\
\texttt{clip\_c2}& $\mathrm{clip}(2y_1,-1,1)+y_2$ & the clipped coordinate \emph{at} the optimum & gain large \\
\midrule
\multicolumn{4}{l}{\emph{Benchmark analogs}}\\
\texttt{sum}     & $y_1+y_2$ & how the sum splits & Truss, RCM40 \\
\texttt{product} & $y_1y_2$ & how it factors, and sign & DTLZ2 \\
\texttt{norm}    & $\sqrt{y_1^2+y_2^2}$ & direction & Color \\
\texttt{max}     & $\max(y_1,y_2)$ & the smaller value & Lake \\
\texttt{ratio}   & $\frac{u_1}{0.05+u_1+u_2}$, $u=\mathrm{sigmoid}(y)$ & scale & RGB \\
\bottomrule
\end{tabular}
\end{center}

\subsection{Measurements}

\begin{center}
\small
\begin{tabular}{lrrrrrl}
\toprule
Map & $\Ulost$ & \multicolumn{3}{c}{observing $h$} & knowing $g$ & Verdict \\
\cmidrule(lr){3-5}
 & & raw & corrected & factor & raw & \\
\midrule
\texttt{identity}      & $0.000$ & $-0.000$ & exact    & $1.0\times$   & $+0.409$ & \textbf{as predicted} \\
\texttt{exp\_b1}       & $0.000$ & $-0.000$ & exact    & $1.0\times$   & $+0.625$ & \textbf{as predicted} \\
\texttt{sinh\_b1}      & $0.000$ & $-0.000$ & exact    & $1.0\times$   & $+0.444$ & \textbf{as predicted} \\
\midrule
\texttt{shiftsq\_c0}   & $0.280$ & $+1.899$ & $+1.490$ & $31\times$    & $+0.115$ & \\
\texttt{shiftsq\_c0.5} & $0.283$ & $+1.688$ & $+1.279$ & $19\times$    & $+0.625$ & sweep runs \\
\texttt{shiftsq\_c1}   & $0.283$ & $+2.303$ & $+1.894$ & $78\times$    & $-0.100$ & \textbf{backwards} \\
\texttt{shiftsq\_c2}   & $0.189$ & $+3.159$ & $+2.750$ & $563\times$   & $-0.726$ & \\
\texttt{shiftsq\_c3}   & $0.232$ & $+2.923$ & $+2.514$ & $326\times$   & $-1.176$ & \\
\midrule
\texttt{abs}           & $0.474$ & $+0.352$ & $-0.057$ & $0.9\times$   & $-0.006$ & unlike \texttt{shiftsq\_c0} \\
\texttt{cos\_w1}       & $2.434$ & $+2.861$ & $+2.452$ & $283\times$   & $-0.793$ & \\
\texttt{cos\_w2}       & $1.284$ & $+1.471$ & $+1.062$ & $12\times$    & $+0.702$ & not monotone in $\omega$ \\
\texttt{cos\_w4}       & $1.161$ & $+1.487$ & $+1.078$ & $12\times$    & $-0.797$ & \\
\texttt{relu}          & $0.075$ & $+0.343$ & $-0.066$ & $0.9\times$   & $+0.201$ & \textbf{as predicted} \\
\texttt{clip\_c2}      & $0.580$ & $+1.517$ & $+1.108$ & $13\times$    & $+0.281$ & \textbf{as predicted} \\
\midrule
\texttt{sum}           & $0.702$ & $+0.229$ & $-0.180$ & $0.7\times$   & $+0.117$ & \textbf{as predicted} \\
\texttt{product}       & $1.640$ & $+0.622$ & $+0.213$ & $1.6\times$   & $-0.411$ & \\
\texttt{norm}          & $0.426$ & $+0.314$ & $-0.095$ & $0.8\times$   & $-0.024$ & \\
\texttt{max}           & $0.358$ & $+0.548$ & $+0.139$ & $1.4\times$   & $+0.067$ & \\
\texttt{ratio}         & $0.267$ & $+0.526$ & $+0.117$ & $1.3\times$   & $-0.418$ & \\
\bottomrule
\end{tabular}
\end{center}

\noindent Paired standard errors run from $0.06$ to $0.48$; the full set is in
\texttt{run.py --collect}. ``Factor'' is $10^{\text{corrected}}$, the ratio of
final regrets.

\subsection{What held}

\paragraph{The invertible controls are exact.} $-0.000\pm0.000$ on all three.
When $g$ is injective the latent arm recovers $h$ from $y$ and becomes the
composite arm, and the code reproduces that identity to machine precision. This
is the check that failed at $+0.956$ before the defects of
Section~\ref{sec:defects} were fixed.

\paragraph{Knowing $g$ is worth nothing on a linear map.} $+0.117\pm0.340$ on
\texttt{sum}, consistent with zero. The prediction was that a linear combination
of independent GPs is itself a GP, so the latent arm collapses to a direct GP on
$y$ with a particular kernel. It does. This is the prediction that bears on the
paper: Truss and RCM40 have linear outer maps, so whatever composite gain they
show cannot come from knowing $g$.

\paragraph{Where the information is destroyed matters more than how much.} This
is the substantive result. \texttt{relu} and \texttt{clip\_c2} both censor, and
\texttt{relu} is the \emph{more} destructive of the two by construction --- it
collapses an entire half-line to a point. Yet:

\begin{center}
\begin{tabular}{llrr}
\toprule
Map & Censoring sits & $\Ulost$ & observing $h$ (corrected) \\
\midrule
\texttt{relu}     & away from the optimum & $0.075$ & $-0.066\pm0.209$ \\
\texttt{clip\_c2} & at the optimum        & $0.580$ & $+1.108\pm0.421$ \\
\bottomrule
\end{tabular}
\end{center}

The rest of the table says the same thing. Sorting by $\Ulost$ does not sort by
gain: \texttt{product} destroys the most information of any analog
($\Ulost=1.64$) and pays $+0.213$; \texttt{sum} destroys $0.702$ and pays
$-0.180$; \texttt{abs} destroys $0.474$ and pays $-0.057$. Whether the ambiguity
falls where the optimizer is looking is what separates them.

If this holds it sharpens the original hypothesis considerably. The quantity to
compute for a new problem is not ``is $g$ invertible'' but ``is $g$ invertible
\emph{near the Pareto front}''.

\subsection{What did not hold: the sweep}
\label{sec:sweepanomaly}

The sweep was designed as the cleanest experiment in the set: one functional
form, one parameter, information loss falling monotonically as $c$ grows. The
gain from observing $h$ should have fallen with it. It rises --- $+1.49$,
$+1.28$, $+1.89$, $+2.75$, $+2.51$ --- and the largest gains occur where $g$ is
\emph{most} nearly injective.

$\Ulost$ explains why: it is flat at $0.19$--$0.28$ across the whole sweep,
where it should fall by two orders of magnitude. The knob does turn on the
initial design --- median $\Ulost$ at iteration $0$ is $1.026$, $0.869$,
$0.104$, $0.013$ for $c=0,1,2,3$, which is exactly the designed behavior --- and
then, over BO iterations, it falls for small $c$ and \emph{rises} for large $c$.
More data cannot increase uncertainty about points already observed.

A second symptom points the same way: ``knowing $g$'' goes negative at $c\ge1$,
reaching $-1.176$. The latent arm holds the generator's true kernel and knows
$g$ exactly; losing to a direct GP by a factor of fifteen in final regret is not
a null result.

Three candidate explanations were tested and two eliminated:

\begin{center}
\begin{tabular}{lll}
\toprule
Hypothesis & Test & Outcome \\
\midrule
Sampler budget too small & $4$ chains / burn $800$ / keep $1500$ vs.\ the loop's
$2/150/200$ & \textbf{ruled out}: $\Ulost$ unchanged \\
Warm starts accumulate error & cold fit vs.\ warm path, identical final design &
\textbf{ruled out}: no consistent direction \\
Level-set moves fail & acceptance in the loop vs.\ in \texttt{--check} &
\textbf{confirmed}, but see below \\
\bottomrule
\end{tabular}
\end{center}

Move acceptance is $0.006$--$0.074$ inside the BO loop and exactly $0.000$ on
clustered designs, against $0.5$--$0.68$ on the spread Sobol designs
\texttt{--check} uses. Single-site branch flips are almost always rejected once
neighbouring points are strongly coupled under the prior, so the sampler stops
traversing branches exactly when the design tightens --- late in the run, when it
matters. The fix is block moves: flip every point within a lengthscale of a
random centre, which keeps the proposal symmetric and the prior ratio in closed
form.

That is a real defect, but it cannot be the whole story, because
\texttt{identity} has no branches to move between and still shows a residual.

\subsection{The residual on \texttt{identity}}
\label{sec:residual}

Run with \texttt{--no-shortcuts}, \texttt{identity} gives
``observing $h$'' $=+0.409\pm0.133$ against a true value of $0$, with
$\Ulost=0.0016$. The sampler recovers $h$; the gap is not information.

The remaining asymmetry is how the predictive is formed. The composite arm draws
$n_{mc}$ samples from one exact Gaussian conditional. The latent arm spreads the
same $n_{mc}$ across a mixture over latent samples, so it covers two sources of
randomness with one budget. Raising $n_{mc}$ from $256$ to $4096$ moves the gap
from $+0.409\pm0.133$ to $+0.275\pm0.168$ --- it shrinks, but it does not vanish.
About $+0.27$ is still unaccounted for on a problem where the latent arm
provably has complete information.

This is the smallest case that still exhibits the defect: one intermediate, no
multimodality, no level-set structure, a posterior that is nearly a point. It is
therefore the right place to debug, and everything in
Section~\ref{sec:sweepanomaly} is downstream of it.

\subsection{The analogs against the benchmark results}

\begin{center}
\begin{tabular}{llrl}
\toprule
Analog & Benchmark & observing $h$ (corrected) & Composite gain in the paper \\
\midrule
\texttt{product} & DTLZ2        & $+0.213$ & largest ($121\%$) \\
\texttt{max}     & Lake         & $+0.139$ & moderate \\
\texttt{ratio}   & RGB          & $+0.117$ & none \\
\texttt{norm}    & Color        & $-0.095$ & moderate \\
\texttt{sum}     & Truss, RCM40 & $-0.180$ & weakest \\
\bottomrule
\end{tabular}
\end{center}

The ordering matches at both ends: the DTLZ2 analog pays the most and the
Truss/RCM40 analog pays the least, which is what the benchmark study found.
Stated carefully, though, these five numbers carry standard errors of $0.06$ to
$0.25$, so no individual pair is separated significantly, and the correction
subtracted from all of them is itself uncertain. Treat this as consistent with
the benchmark results rather than as confirmation of them.

\subsection{Predictions this generates}

\begin{enumerate}
\item \textbf{Local, not global, invertibility should predict the composite
gain.} Compute $\Ulost$ on each real benchmark restricted to near-Pareto-optimal
points, and it should correlate with the measured composite gain better than
$\Ulost$ computed over the whole domain. The \texttt{relu}/\texttt{clip}
contrast predicts this directly and it is cheap to test --- no new BO runs, just
a posterior fit on the points each existing run already evaluated.

\item \textbf{Truss and RCM40 should gain from the extra observations only.}
Their outer maps are linear, so knowing $g$ is worth nothing; the whole gain
must be the $k$ numbers per evaluation. A direct test is to give the direct arm
$k$ evaluations per iteration instead of one --- if the gain disappears, it was
never about composite structure at all.

\item \textbf{A kink at the optimum should not matter, but ambiguity should.}
\texttt{abs} and \texttt{shiftsq\_c0} destroy the same information and differ
only by smoothness; they measure $-0.057$ and $+1.490$. If that survives the
defects above it says the acquisition function's difficulty with a kink is a
larger effect than the inference problem --- worth isolating on its own.

\item \textbf{The number of preimages should not be the controlling variable.}
$\cos(\omega z)$ at $\omega=1,2,4$ gives $+2.45$, $+1.06$, $+1.08$ --- the
\emph{fewest} preimages pays the most. If this holds it contradicts the natural
reading of ``how much information is destroyed'' and supports the same local
story as the censoring pair.
\end{enumerate}

% ======================================================================
\part{Limitations}
\label{part:limits}
% ======================================================================

\begin{itemize}
\item \textbf{The hyperparameter asymmetry is not yet removed}, only measurable.
See Section~\ref{sec:hyperconfound}. ``Knowing $g$'' is an upper bound until the
all-MLE run is done.

\item \textbf{Noise makes injectivity a matter of degree.} As measured in
Check~2, an injective $g$ still loses information where its slope is small
relative to $\sigma$. The clean statement holds only in the noise-free limit.

\item \textbf{Continuous fibers mix more slowly than discrete ones.} Reflections
for $z^2$ are single moves that jump between modes; following the circle fiber
of a norm is a single-site random walk over a smooth field, which mixes more
slowly when observed points are strongly coupled. The ratio map has no exact
moves at all and relies on elliptical slice sampling alone, so it is the least
trustworthy condition in the set.

\item \textbf{$\Ulost$ is a variance, not an information measure.} The
non-monotonicity at $c=0.5$ shows the two can come apart when posterior modes are
unequally weighted.

\item \textbf{Nothing here has $k>2$.} The real benchmarks have four to eight
intermediates. Whether the latent arm's advantage scales with the number of
hidden dimensions is untested, and it is also where the sampler would be least
reliable, so this is a deliberate omission rather than an oversight.

\item \textbf{The reference minimum is estimated}, from a large Sobol set refined
by gradient ascent, so regret is measured against a very good estimate rather
than a proven optimum.

\item \textbf{Single objective.} The decomposition is cleanest with one
objective; whether the same accounting holds once a hypervolume acquisition is
involved is a separate question, and the natural next step once these results are
in.
\end{itemize}

\appendix
\section{Where each piece lives in the code}

\begin{center}
\begin{tabular}{lll}
\toprule
Design element & File & Object \\
\midrule
Random Fourier feature draws, standardization & \texttt{test\_functions.py} & \texttt{RFFDraw} \\
The nineteen outer maps & \texttt{test\_functions.py} & \texttt{OUTER\_MAPS} \\
Level-set moves per map & \texttt{test\_functions.py} & \texttt{OuterMap.moves} \\
Named groups for submission & \texttt{test\_functions.py} & \texttt{GROUPS} \\
Shared prior for arms 2 and 3 & \texttt{arms.py} & \texttt{KnownPrior} \\
Arm 1 & \texttt{arms.py} & \texttt{DirectArm} \\
Arm 2 & \texttt{arms.py} & \texttt{CompositeArm} \\
Arm 3, sampler and shortcuts & \texttt{arms.py} & \texttt{LatentArm} \\
Elliptical slice step & \texttt{arms.py} & \texttt{LatentArm.\_ess\_step} \\
Level-set Metropolis moves, Eq.~\eqref{eq:priorratio} & \texttt{arms.py} & \texttt{LatentArm.\_level\_set\_moves} \\
Projection onto $\{g(h)=y\}$ & \texttt{arms.py} & \texttt{LatentArm.\_project} \\
Exact Gaussian posterior for linear $g$ & \texttt{arms.py} & \texttt{exact\_linear\_posterior} \\
Acquisition & \texttt{run.py} & \texttt{log\_ei}, \texttt{optimize\_acq} \\
Paired differences and standard errors & \texttt{run.py} & \texttt{paired} \\
The three validation checks & \texttt{run.py} & \texttt{run\_checks} \\
Convergence diagnostics & \texttt{run.py} & \texttt{\_mixing} \\
Pooling over function draws & \texttt{run.py} & \texttt{main}, \texttt{--collect} \\
Cluster submission & \texttt{cluster/} & \texttt{three\_arm.sub}, \texttt{submit\_three\_arm.sh} \\
\bottomrule
\end{tabular}
\end{center}

\begin{thebibliography}{9}
\bibitem{murray2010}
I.~Murray, R.~P. Adams, and D.~J.~C. MacKay.
Elliptical slice sampling.
In \emph{Proceedings of the 13th International Conference on Artificial
Intelligence and Statistics}, PMLR 9:541--548, 2010.
\end{thebibliography}

\end{document}
